\honest{Practical Advice Backed By Deep Theories}{Eliezer Yudkowsky}{2009}

Seth Roberts lost weight on holiday in Europe while drinking unfamiliar-tasting drinks. Had he known nothing of the research on set points and flavor-calorie learning, he would have written one blog post about amazing fruit juices, the trick would have worked briefly for a few people, and there would have been no Shangri-La Diet.

The real diet does not work for everyone; it does not work for me, and there is ``no apparent reason for this.'' But ``the reason why as many people have benefited as they have'' is that Roberts knew the science. \nb{How many benefited was not measured. Fifteen days earlier, in ``The Unfinished Mystery of the Shangri-La Diet,'' Yudkowsky had written that the theory ``has never been analyzed by controlled experiment,'' and had cited reports from forums and bloggers.}

My own advice worked the same way. Readers named ``the bottom line'' and ``engines of cognition'' as the most important things they learned here. \nb{In the thread the post links, three of 99 comments name each. That shows which ideas readers valued, not that the theory changed what they did.} Telling you to keep an open mind and follow the evidence would have been less useful without probability theory, which says what counts as evidence.

Causal theories from textbooks are likelier to be true than ones made up on the spot, and theory-backed advice is ``the brand that can distinguish LW from ten thousand other blogs.''

If I tell you that satisfaction with food probably depends more on quality than on quantity, you will forget it. If I tell you about scope insensitivity, duration neglect and the peak-end rule, you will see at your plate that a large portion leaves almost the same memory as a small one. ``(You also know to save the dessert for last.)'' \nb{A study of meals published in 2007, by Rode, Rozin and Durlach, supports the point about portions: ``small portions of a highly favored dish will have roughly the same memorial effect as large portions.'' The same study found ``no evidence for peak, primacy or recency effects,'' so the dessert tip was not borne out for meals in that test.}

I want advice on akrasia from someone who has read the experimental psychology of willpower, beaten their own akrasia, and can explain how in terms of replicable experiments. Most advisers, ``for the most part,'' describe mental levers they cannot point to. \nb{Such work existed: Piers Steel's 2007 meta-analysis tied procrastination to hyperbolic discounting, and a 2011 LessWrong post by Luke Muehlhauser built its advice on it.}

Backing gets harder to cite well as it goes from experimental results (``p < 0.05 may fail to replicate''), to true causal accounts, to mathematics validly interpreted. If you do not know whom to trust, start with experiments. Practical advice ``really, really'' does become more powerful with all three. \nb{The post supports this with a counterfactual, readers' favorite posts and an imagined dinner plate. It reports no comparison of advice given with and without a theory.}
